% Methodology section: Uncertainty Quantification for pediatric bone age regression.
% Self-contained, compilable document. Can also be \input{} into a larger report
% by reusing the body between \begin{document} and \end{document}.
\documentclass[11pt]{article}

\usepackage[utf8]{inputenc}
\usepackage[T1]{fontenc}
\usepackage{amsmath}
\usepackage{amssymb}
\usepackage{graphicx}
\usepackage{booktabs}
\usepackage{geometry}
\usepackage[hidelinks]{hyperref}

\geometry{margin=1in}

\newcommand{\code}[1]{\texttt{#1}}

\title{Uncertainty Quantification for Pediatric Bone Age Regression:\\ Methodology}
\author{}
\date{}

\begin{document}

\maketitle

\section{Overview and Design Philosophy}

This work augments a deterministic point-regression model for pediatric bone age
estimation with three complementary uncertainty quantification (UQ) methods:
\emph{Monte Carlo (MC) Dropout}, \emph{Heteroscedastic Regression}, and a
\emph{Bayesian Neural Network (BNN)}. Each method produces, for every radiograph,
not only a point estimate of bone age (in months) but also a calibrated
\emph{prediction interval} expressing the model's confidence in that estimate.

The UQ layer is deliberately designed to be \emph{additive}. The production
point-regression pipeline---the model definition (\code{model.py}), the training
loop (\code{train.py}), and the data pipeline (\code{data\_loader.py})---is left
unchanged. Every UQ method lives in its own self-contained package under
\code{UQ/<method>/} and reuses a small set of shared utilities under
\code{UQ/common/}. This separation guarantees that (i) the deterministic baseline
remains reproducible, (ii) a method can be removed simply by deleting its folder,
and (iii) all methods are evaluated through identical interval-construction,
calibration, and metric code, so they can be ranked fairly.

A unifying principle across the three methods is that they all reduce, at
inference time, to a per-sample predictive mean $\mu_i$ and a per-sample
predictive standard deviation $\sigma_i$. From these two quantities, Gaussian
prediction intervals are constructed and scored on the same coverage and width
metrics. All metrics are computed in de-normalized \emph{month space} so that
the methods are directly comparable to one another and to the deterministic
baseline.

\section{Shared Experimental Setup}
\label{sec:setup}

\subsection{Base Architecture}

All three UQ methods share the same backbone as the deterministic baseline, a
multi-input convolutional regression network (\code{MultiInputModel} in
\code{model.py}). The architecture fuses image features with patient sex:

\begin{itemize}
  \item \textbf{Image branch.} A pre-trained EfficientNet-B3 (ImageNet
  \code{IMAGENET1K\_V1} weights) acts as the feature extractor. Its
  classification head is replaced by an identity mapping
  (\code{nn.Identity}), exposing the raw pooled feature vector of dimension
  $d_\text{img}$ (the EfficientNet-B3 classifier input width).
  \item \textbf{Sex branch.} The binary sex indicator is passed through a linear
  layer $\mathbb{R}^{1}\!\to\!\mathbb{R}^{32}$ followed by a ReLU activation.
  \item \textbf{Fusion head.} The image and sex features are concatenated into a
  $(d_\text{img}+32)$-dimensional vector and passed through a fully connected
  layer to $256$ units (ReLU), a dropout layer with rate $0.5$, and a final
  linear layer producing a single scalar---the normalized bone age.
\end{itemize}

Each UQ method modifies \emph{only} the prediction head of this backbone (or its
inference procedure), as detailed in Sections~\ref{sec:mcdropout}--\ref{sec:bnn}.

\subsection{Data, Preprocessing, and Splits}

Radiographs are loaded as RGB images, resized to $300\times300$ pixels, and
normalized with the standard ImageNet channel statistics
($\mu=[0.485, 0.456, 0.406]$, $\sigma=[0.229, 0.224, 0.225]$). During training,
the images are augmented with random rotation ($\pm20^\circ$), random horizontal
flips, random affine translation/scaling, and color jitter
(brightness/contrast). Evaluation and calibration use the deterministic
resize-and-normalize transform with no augmentation and no shuffling, so that
predictions remain aligned with their targets.

The training source (\code{train.csv}) is partitioned, stratified by sex with a
fixed random seed (\code{RANDOM\_STATE}~$=42$), into three disjoint splits in a
$50/25/25$ ratio:
\begin{itemize}
  \item \textbf{train} ($50\%$): used to fit model parameters;
  \item \textbf{val} ($25\%$): used to monitor point accuracy and drive the
  learning-rate scheduler;
  \item \textbf{calib} ($25\%$): a held-out split reserved for \emph{post-hoc
  uncertainty calibration} (Section~\ref{sec:calibration}).
\end{itemize}
The RSNA held-out \code{test.csv} (images under \code{data/test/}) serves as the
external test set and is never seen during training or calibration.

Bone age targets $y$ (in months) are normalized to $\tilde{y} = y / y_{\max}$,
where $y_{\max}$ is the maximum bone age in the \emph{training} split only;
applying the train-derived $y_{\max}$ to every split avoids information leakage.
All model outputs are produced in this normalized space and multiplied back by
$y_{\max}$ (\emph{de-normalization}) before any metric is computed, so that
coverage and interval widths are reported in clinically interpretable months.

The dedicated \code{calib} split is what makes principled post-hoc calibration
possible: it provides labelled data, disjoint from both training and test, on
which the relationship between predicted uncertainty and observed error can be
estimated.

\section{Shared Uncertainty Infrastructure}
\label{sec:infra}

The three methods share the interval-construction, metric, and calibration code
under \code{UQ/common/}. Describing it once avoids repetition and makes explicit
that differences in results stem from the methods themselves, not from
differing evaluation code.

\subsection{Gaussian Prediction Intervals}

Given a per-sample predictive mean $\mu_i$ and standard deviation $\sigma_i$ (both
in months), a two-sided prediction interval at confidence level $c$ is
constructed under a Gaussian assumption as
\begin{equation}
  \big[\,\mu_i - z_c\,\sigma_i,\;\; \mu_i + z_c\,\sigma_i\,\big],
  \label{eq:interval}
\end{equation}
where $z_c$ is the standard-normal quantile for the nominal level: $z_{90}=1.645$
for $90\%$ and $z_{95}=1.960$ for $95\%$. The two confidence levels share the
same $\mu_i$ and $\sigma_i$ and differ only through $z_c$.

\subsection{Evaluation Metrics}
\label{sec:metrics}

Let $\{(y_i, \mu_i)\}_{i=1}^{N}$ be the true ages and predictive means, and let
$[\ell_i^{(c)}, u_i^{(c)}]$ denote the interval at level $c$. Two
interval-quality metrics are reported at both $90\%$ and $95\%$:

\paragraph{Prediction Interval Coverage Probability (PICP).} The empirical
fraction of true values contained in their intervals,
\begin{equation}
  \mathrm{PICP}@c = \frac{1}{N}\sum_{i=1}^{N}
  \mathbb{1}\!\left[\,\ell_i^{(c)} \le y_i \le u_i^{(c)}\,\right].
\end{equation}
A well-calibrated $90\%$ interval should yield $\mathrm{PICP}@90 \approx 0.90$.

\paragraph{Mean Prediction Interval Width (MPIW).} The average interval width in
months,
\begin{equation}
  \mathrm{MPIW}@c = \frac{1}{N}\sum_{i=1}^{N}
  \left(u_i^{(c)} - \ell_i^{(c)}\right).
\end{equation}
Given adequate coverage, smaller MPIW is preferable, as it indicates tighter,
more informative intervals.

\paragraph{Point-accuracy metrics.} For parity with the deterministic baseline,
the mean prediction $\mu_i$ is also scored with mean absolute error (MAE),
root-mean-square error (RMSE), the coefficient of determination $R^2$, and the
fraction of predictions within $6$ and $12$ months of ground truth. A method is
considered well calibrated when its PICP lands close to the nominal level; among
methods with comparable coverage, the one with smaller MPIW is preferred.

\subsection{Post-hoc Calibration via Temperature Scaling}
\label{sec:calibration}

The raw predictive standard deviation produced by a UQ method is rarely
perfectly scaled: in practice it tends to be too small, so the intervals of
Equation~\eqref{eq:interval} \emph{under-cover} (PICP below nominal). The
standard remedy is a single-parameter \emph{temperature scaling} of the
standard deviation. A scalar $s>0$ is fit on the held-out \code{calib} split as
the maximum-likelihood scale of a zero-mean Gaussian fit to the standardized
residuals:
\begin{equation}
  s = \sqrt{\frac{1}{N}\sum_{i=1}^{N}\left(\frac{y_i - \mu_i}{\sigma_i}\right)^{2}},
  \label{eq:temperature}
\end{equation}
computed in month space, with $\sigma_i$ floored by a small $\epsilon$ to avoid
division by zero. The calibrated intervals then use $\sigma_i^{\text{eff}} = s\,
\sigma_i$ in Equation~\eqref{eq:interval}. Intuitively, $s=1$ means the raw
standard deviation already matched the residual spread; $s>1$ widens the
intervals (the common case), and $s<1$ tightens them. Because a single scalar is
applied across both confidence levels, the $90\%$ and $95\%$ intervals remain
mutually consistent.

The fitted $s$ is persisted to a small, self-describing JSON file (recording
the scale, the split it was fit on, the sample count, a timestamp, and method
metadata) so it can be audited and reused. The same $s$ that is estimated on
\code{calib} is later applied unchanged to the \code{test} split, which is the
correct protocol: the test set is used only for final measurement, never for
fitting calibration.

\subsection{Common Workflow}

Every method follows the same three-stage protocol:
\begin{enumerate}
  \item \textbf{Train} the model and select a checkpoint (selection criteria
  described per method below);
  \item \textbf{Fit calibration} by running inference on the \code{calib} split
  and computing $s$ via Equation~\eqref{eq:temperature};
  \item \textbf{Evaluate} on the \code{test} split with the saved $s$ applied,
  emitting per-sample predictions and a single metrics row.
\end{enumerate}

\section{Method 1: Monte Carlo Dropout}
\label{sec:mcdropout}

\subsection{Principle}

Monte Carlo Dropout treats the dropout layer already present in the trained
network as an approximate source of epistemic (model) uncertainty. Whereas a
standard model disables dropout at inference (\code{model.eval()}), MC Dropout
keeps the dropout layers \emph{active} while leaving all other layers (e.g.\
BatchNorm) in evaluation mode. Each forward pass then samples a different
sub-network, and repeated passes over the same input yield a distribution of
predictions.

Concretely, the model is put into evaluation mode and every \code{nn.Dropout}
module is switched back to training mode. For an input $x_i$, $T$ stochastic
forward passes produce predictions $\{\hat{y}_i^{(t)}\}_{t=1}^{T}$ (in normalized
space). The predictive mean and standard deviation are the per-sample sample
statistics
\begin{equation}
  \mu_i = \frac{1}{T}\sum_{t=1}^{T}\hat{y}_i^{(t)},
  \qquad
  \sigma_i = \sqrt{\frac{1}{T}\sum_{t=1}^{T}\big(\hat{y}_i^{(t)} - \mu_i\big)^{2}},
\end{equation}
which are de-normalized to months before interval construction. The default
number of passes is $T=30$ at evaluation time. The data loader does not shuffle,
so targets stay aligned across passes.

A notable property of this setup is that the baseline architecture contains a
\emph{single} dropout layer in its head (the EfficientNet classifier having been
replaced by an identity mapping). The MC variance is therefore driven entirely
by that one layer. The model is intentionally left unchanged so that the
``same regression model'' comparison against the other methods remains valid.

\subsection{UQ-aware Training and Checkpoint Selection}

Post-hoc calibration (Section~\ref{sec:calibration}) can pin coverage to the
nominal level, but it cannot make the intervals \emph{tighter}; interval
sharpness is determined at training time. A dedicated training loop
(\code{train\_mc\_dropout.py}), a sibling of the base trainer, trains the same
architecture while optimizing for uncertainty quality:
\begin{itemize}
  \item The dropout rate is configurable (rather than hard-coded), since it
  governs the magnitude of the MC predictive spread.
  \item After each epoch, MC Dropout is run on the \code{calib} split, a
  temperature $s$ is fit, and the epoch is scored by the quality of the
  \emph{calibrated} intervals rather than by plain validation MAE.
\end{itemize}

The checkpoint-selection score (lower is better) can be one of:
\begin{itemize}
  \item \code{val\_mae}: plain validation MAE (the baseline behaviour);
  \item \code{picp90}: $\lvert \mathrm{PICP}@90 - 0.90 \rvert$, i.e.\ closeness
  of calibrated coverage to nominal;
  \item \code{mpiw90}: calibrated $\mathrm{MPIW}@90$ (tightest intervals);
  \item \code{combo} (default): a trade-off
  $\mathrm{MPIW}@90 + \lambda\,\mathrm{MAE}$, where $\lambda$ (default $0.01$)
  weights the calibration-split MAE in months.
\end{itemize}
The rationale for \code{combo} is that, since calibration already fixes
coverage, the meaningful objective becomes minimizing interval width while a
small MAE penalty prevents point accuracy from regressing. To keep training
affordable, the per-epoch calibration evaluation uses a small number of passes
(default $10$), while the final, definitive calibration after restoring the best
checkpoint uses more passes (default $30$). The resulting temperature is saved
for reuse at test time.

\section{Method 2: Heteroscedastic Regression}
\label{sec:hetero}

\subsection{Principle}

Heteroscedastic regression keeps the same EfficientNet-B3 + sex backbone but
replaces the single regression head with \emph{two} heads sharing the
$256$-dimensional fused representation:
\begin{itemize}
  \item \code{mean\_out} predicts the normalized bone age $\mu(x)$ (the same
  target as the baseline);
  \item \code{log\_var\_out} predicts the logarithm of the per-sample variance,
  $\log \sigma^2(x)$, of a Gaussian observation model.
\end{itemize}
Predicting $\log \sigma^2$ rather than $\sigma$ or $\sigma^2$ directly keeps the
output unconstrained over $(-\infty, \infty)$, avoiding negative variances and
improving numerical stability. The predictive standard deviation is recovered
as $\sigma(x) = \exp\!\big(\tfrac{1}{2}\log\sigma^2(x)\big)$.

This formulation captures \emph{input-dependent aleatoric uncertainty}: the
network can assign wider intervals to harder radiographs and tighter ones to
easier cases. Unlike MC Dropout, inference requires only a \emph{single}
deterministic forward pass that returns $(\mu, \log\sigma^2)$ jointly, making it
considerably cheaper at evaluation time. It is worth noting that this method
models aleatoric (data) uncertainty only and does not, on its own, capture
epistemic uncertainty about the model weights; in this respect it is
complementary to MC Dropout and the BNN.

\subsection{Training Objective}

The model is trained by minimizing the Gaussian negative log-likelihood (NLL).
For a per-sample Gaussian $\mathcal{N}\!\big(\mu(x), \sigma^2(x)\big)$, the NLL,
up to an additive constant, is
\begin{equation}
  \mathcal{L}_\text{NLL} =
  \frac{1}{2}\,\Big(\exp\!\big(-\log\sigma^2\big)\,(y - \mu)^2 + \log\sigma^2\Big).
  \label{eq:nll}
\end{equation}
The two terms of Equation~\eqref{eq:nll} balance one another. The first term
divides the squared residual by the predicted variance, so that the model can
attenuate the loss contribution of inputs it deems uncertain by predicting a
larger variance. The second term, $\log\sigma^2$, is a penalty that prevents the
trivial solution of predicting infinite variance everywhere.

To prevent numerical pathologies, the predicted log-variance is clamped to
$[-12, 4]$ (corresponding to $\sigma$ in roughly $[1.8\times10^{-3}, 7.4]$ in
normalized age space). This guards against two failure modes: a log-variance
that is too negative makes the inverse variance $\exp(-\log\sigma^2)$ explode
(causing exploding gradients), while one that is too positive lets training
collapse into predicting massive uncertainty.

An optional auxiliary term may be mixed in: a SmoothL1 loss on the mean head,
weighted by $\lambda_\text{mae}$,
\begin{equation}
  \mathcal{L} = \mathcal{L}_\text{NLL}
  + \lambda_\text{mae}\,\mathrm{SmoothL1}(\mu, y).
\end{equation}
Setting $\lambda_\text{mae}=0$ recovers pure NLL; a small positive value (e.g.\
$0.01$) keeps the mean head close to a SmoothL1 optimum, which empirically
preserves point accuracy that pure NLL might otherwise sacrifice in favour of
better-calibrated variance.

\subsection{Inference and Selection}

At inference, a single forward pass yields $(\mu, \log\sigma^2)$; the clamped
log-variance is converted to $\sigma$, both are de-normalized, the temperature
$s$ from Section~\ref{sec:calibration} is applied, and intervals are formed via
Equation~\eqref{eq:interval}. Checkpoint selection mirrors the MC Dropout
trainer, supporting \code{val\_mae} (the default), \code{picp90}, \code{mpiw90},
and \code{combo}; when a UQ-aware criterion is selected, the calibrated interval
quality is evaluated on the \code{calib} split. The final temperature is fit on
\code{calib} after restoring the best checkpoint and saved for the test-time
evaluation.

\section{Method 3: Bayesian Neural Network}
\label{sec:bnn}

\subsection{Principle}

The Bayesian Neural Network captures \emph{epistemic} uncertainty by treating
the weights of the regression head as random variables with a learned posterior
distribution, following the Bayes-by-Backprop framework of
Blundell et al.~(2015). The two fully connected head layers (\code{fc1} and
\code{out}) are replaced by \code{VariationalLinear} layers, while the
pre-trained EfficientNet-B3 backbone remains deterministic. Restricting the
Bayesian treatment to the head keeps training tractable and stable while still
modelling weight uncertainty where the regression decision is made.

\subsection{Variational Layers}

Each \code{VariationalLinear} layer maintains, for every weight and bias, a mean
parameter $\mu$ and an unconstrained scale parameter $\rho$. The positive
standard deviation is obtained through a softplus transform,
$\sigma = \mathrm{softplus}(\rho) = \log\!\big(1 + e^{\rho}\big)$, which keeps
$\sigma > 0$ without explicit constraints. During a stochastic forward pass,
weights are drawn from the Gaussian posterior using the reparameterization
trick,
\begin{equation}
  w = \mu + \sigma \odot \varepsilon, \qquad \varepsilon \sim \mathcal{N}(0, I),
\end{equation}
which keeps the sampling operation differentiable with respect to $\mu$ and
$\rho$. The means $\mu$ are initialized with He-uniform initialization, and the
scales are initialized so that the initial posterior standard deviation is small
($\sigma \approx 0.05$, via $\rho = \log(e^{\sigma}-1)$). At a deterministic
(mean) forward pass the layer simply uses $w=\mu$.

\subsection{Training Objective (ELBO)}

Training maximizes the Evidence Lower Bound (ELBO), expressed here as a loss to
be minimized:
\begin{equation}
  \mathcal{L}_\text{ELBO} =
  \mathrm{SmoothL1}(y, \hat{y})
  + \beta\,\frac{\mathrm{KL}}{N_\text{train}},
  \label{eq:elbo}
\end{equation}
where the SmoothL1 term is the (negative log-likelihood surrogate) data-fit
term, $N_\text{train}$ is the number of training samples, and $\beta$ is a
configurable weight (\code{kl\_weight}) on the Kullback--Leibler regularizer.
Dividing the KL term by $N_\text{train}$ scales the per-dataset complexity cost
to a per-sample basis, balancing it against the per-sample data term.

The KL divergence is accumulated analytically over all variational weights and
biases. For a Gaussian posterior $\mathcal{N}(\mu, \sigma^2)$ against a
zero-mean Gaussian prior $\mathcal{N}(0, \sigma_\text{prior}^2)$, each parameter
contributes
\begin{equation}
  D_\text{KL}\!\big(\mathcal{N}(\mu,\sigma^2)\,\|\,\mathcal{N}(0,\sigma_\text{prior}^2)\big)
  = \log\frac{\sigma_\text{prior}}{\sigma}
  + \frac{\sigma^2 + \mu^2}{2\,\sigma_\text{prior}^2}
  - \frac{1}{2},
\end{equation}
where the prior standard deviation $\sigma_\text{prior}$ (\code{prior\_sigma}) is
a hyperparameter.

\subsection{Inference}

At inference, the variational layers (and the head dropout) are kept active so
that each forward pass samples a different set of head weights. As in MC
Dropout, $T$ stochastic passes over an input yield a predictive sample whose
per-sample mean and standard deviation give $\mu_i$ and $\sigma_i$. These are
de-normalized, scaled by the calibration temperature $s$, and turned into
intervals via Equation~\eqref{eq:interval}. The UQ-aware training loop
(\code{train\_bnn.py}) selects checkpoints with the same criteria as the other
methods (defaulting to \code{combo}), evaluating calibrated interval quality on
the \code{calib} split each epoch and fitting the definitive temperature with
more passes after restoring the best checkpoint.

\section{Cross-method Comparison Protocol}
\label{sec:compare}

Because all three methods reduce to a per-sample $(\mu_i, \sigma_i)$ and are
scored by the same code, their results are directly comparable. Every method
emits a metrics row following a shared schema---method name, split, the point
metrics (MAE, RMSE, MSE, $R^2$, fraction within $6$/$12$ months), the interval
metrics (PICP@90/95, MPIW@90/95), the sample count, optional method-specific
fields (e.g.\ number of MC samples, applied temperature, calibration flag), and
a timestamp.

An aggregation step (\code{compare.py}) scans every method's metrics files,
retains the most recent row per (method, split) pair, and ranks the methods.
The ranking key is, first, the absolute calibration error
$\lvert \mathrm{PICP}@90 - 0.90 \rvert$ (best-calibrated first) and, second, the
calibrated $\mathrm{MPIW}@90$ (tightest intervals as a tie-breaker). This
encodes the guiding principle that coverage close to nominal is the primary
requirement, and that among adequately calibrated methods the tighter,
more informative intervals are preferred. The numerical outcomes of this
comparison are reported in the Results section.

\end{document}
